\section{Finite cocycles and the two type-prime local bounds}

% BEGIN T84L1
\paragraph{T84L1. The lifting exponent on a family over a degree-one place.}
Put $p=\ell_i$, $N=n_V+1$, $a=\eta+t+1$. The degree-one test on
the fixed completion identifies $K=E_w$ with $\mathbb Q_p$ inside
the given algebraic closure. Suppose $a<H$ and
$C_T(HN^2+N^3)<p$. In $\mathrm{PotFam}$ take its common finite
normal extension $K'/K$, its Frobenius $F$ and its $F$-stable
set $\Sigma$ of at most $N$ inertia characters containing $1$.
The crystalline period action on the unramified constants of $K'$
is the one in that datum. In $\mathrm{JumpFam}$ take the common
set $J\subset[-a,H]$, containing $0$ and of size at most $N+1$.

The lattice maps tested by $\mathrm{LiftFamW}$ can have arbitrarily
large source and target. Their rational spaces embed equivariantly
in finite sums of family members. After restriction to $K'$ each
is crystalline by the subobject and direct-sum properties of the
Fontaine comparison. Their jump sets at $K$ stay in $J$: tensor
an injection into the sum with $B_{\rm dR}$, and choose an
ordinary linear left inverse of the injection. The induced map
on the period tensor spaces then pulls the degree-$d$ filtration
back to the degree-$d$ filtration for every $d$. Taking Galois
invariants preserves this intersection. A nonzero graded
piece thus injects into the corresponding graded piece of the sum.
Similarly on $D_{\rm cris,K'}$ the injection respects inertia.
Each character in $\Sigma$ is trivial on the indicated subgroup
over $K'$ and has order prime to $p$, by $\mathrm{InertiaThrough}$
at any one family member. On the ambient period space the
action factors through a finite abelian quotient, the image of
these characters. Its character projectors preserve the image
of the injected period space, which is inertia stable by
normality. Thus the same $\Sigma$ works for each tested
subrepresentation. For an empty family all tested spaces are
zero and the lifting assertion is immediate.

Let $r'$ be the largest Frobenius orbit length of
$\alpha\beta^{-1}$ for $\alpha,\beta\in\Sigma$. Frobenius
permutes $\Sigma$ and fixes the trivial character. For two
characters in the same orbit, its length bounds that of their
ratio, and is at most $N-1$ unless the orbit is the trivial
singleton. For two distinct orbits of lengths $d_1,d_2$, the
ratio is fixed after $\operatorname{lcm}(d_1,d_2)$ steps, and
$d_1d_2\le(d_1+d_2)^2/4\le N^2/4$. Since $N\ge2$, this
gives $r'\le\lfloor N^2/4\rfloor$ in all cases. Now apply
the stated $\mathrm{admT}(C_T)$, the quantitative potentially
crystalline lifting theorem, to the two tested lattices and
their surjection. Its hypotheses on these very spaces have
just been verified, including the full labelled jump set.
It gives lifting with exponent $C_T(Nr'+a(N+r'))$. Multiply
the finite lift and its integral coboundary by
$p^{e_{\max}-C_T(Nr'+a(N+r'))}$, where
\[
 e_{\max}=C_T\bigl(N\lfloor N^2/4\rfloor+
                         a(N+\lfloor N^2/4\rfloor)\bigr).
\]
This proves $\mathrm{StepLift}$ with its original constant.
The quantitative lifting theorem is an explicit input here.
% END T84L1

% BEGIN T84L2
\paragraph{T84L2. A cut cocycle as the pushout of an integral rational coboundary.}
Write $L$ for the local coefficient field and $O$ for its
integers. Let $M$ have rank $k$ and satisfy $\mathrm{SubTw}$.
Thus $M$ is continuous and its rational realization embeds
globally in one of the retained twists $r_x\chi_1^{-1}$.
Write $U$ for its local rational space and $T$ for its standard
lattice. A generator of $\mathrm{cutCoc}$ is given, over the
original coefficients, by a finite tuple of retained twists,
a finitely generated stable module $Y$ in their sum, a vector
$E$ with $y(g)=gE-E\in Y$, and an equivariant integral linear
map $f:Y\to U/T$ under the coefficient identifications.

There is a global stable complement to the rational span $X$ of
$Y$. Indeed each good rational representation is semisimple,
twisting by a nonzero character preserves every invariant
subspace and its complements, and a finite direct sum of
semisimple modules is semisimple. Apply the equivariant projection
onto $X$ to $E$. Then $y(g)=gE_X-E_X$ since the projection
fixes every $y(g)$. Transport this decomposition to the
completion and restrict to $G_K$. The actual image of $Y$ under
the given coefficient isomorphism is a stable full lattice in
$X$, being finitely generated, torsion-free and spanning. It
retains the map $f$. Thus $y$ is a continuous integral cocycle
whose rational class in $X$ is zero. In particular it is in the
integral Bloch--Kato finite condition. Its pushout $f(y)$ is
the original local function on $G_K$ even when $E_X\notin Y$.

The same global complement construction applies to $U$ inside
its chosen member. It follows that both $X$ and $U$ embed in
finite local direct sums of the twisted family, and that their
Tate duals $X'=X^*(1)$ and $U'=U^*(1)$ embed in corresponding
sums of the dual family: extend functionals by zero on the
transported complements. The pairings transform by the
cyclotomic character. In either direction all the objects in
the graph construction below therefore pass the rational
embedding tests of $\mathrm{LiftFamW}$. This argument uses
semisimplicity before restriction to a decomposition group.
The common $K'$ of the direct potential family makes $X,U$
and their Tate duals crystalline over $K'$ by Fontaine's
subobject, sum and dual properties. Potentially crystalline
implies de Rham over $K$.
% END T84L2

% BEGIN T84L3
\paragraph{T84L3. The direct and transposed lattices.}
For the direct lifting condition at exponent $e$, choose an
$L$-linear lift $s:X\to U$ of $f$ on $Y$, by lifting the
images of an $O$-basis. In $X\oplus U$ take the lattice
\[
          B=\{(v,u):v\in Y,\ u-sv\in T\}.
\]
A basis is obtained from the graph of the chosen basis of $Y$
together with a basis of $T$. Equivariance of $f$ says exactly
that $gsv-sgv\in T$ for $v\in Y$, so $B$ is stable.
Projection onto $Y$ is surjective. Apply the finite-cocycle
lifting condition of the direct family to $y$. We obtain a
finite $B$-valued cocycle $(b_1,b_2)$ and $v\in Y$ with
$p^ey=b_1+\delta v$, where $\delta v(g)=gv-v$.
Thus $p^ef(y)=\overline{b_2}+\delta f(v)$. Projection takes a
rationally finite cocycle to a finite one, so this lies in
$H_{f,D}(T)$, the reduction of rational finite cocycles plus
all coboundaries with values in $U/T$.

If only the dual family lifts at $e$, put
$Y'=Y^*(1)$, $T'=T^*(1)$ and let $s^*:U'\to X'$ be the
transpose under the perfect pairings with values in $L(1)$.
In $U'\oplus X'$ use
\[
          B'=\{(b,c):b\in T',\ c-s^*b\in Y'\}.
\]
This projects onto $T'$ and is stable. Pair the transpose's
equivariance defect with $v\in Y$ and use the preceding
defect condition for $g^{-1}v$; the cyclotomic factor
preserves $O(1)$. Given any integral finite test cocycle
$d$ with values in $T'$, dual lifting gives a finite
$B'$-valued cocycle $(b,c)$ with $p^ed=b+\delta t'$ for
some $t'\in T'$. In the Tate cup of $b$ with $f(y)$, the
representative at $(g,h)$ is
$\langle g s(y(h)),b(g)\rangle$ modulo $O(1)$.
Its difference from $\langle g y(h),c(g)\rangle$ is integral:
the first error pairs $g s(y(h))-s(g y(h))\in T$ with $b(g)\in T'$,
and the second pairs $g y(h)\in Y$ with $c(g)-s^*b(g)\in Y'$.
The latter rational cup is a coboundary since $y=\delta E_X$
rationally. Reducing its continuous rational primitive modulo
$O(1)$ preserves that identity. Hence the Tate pairing of
$p^ef(y)$ with $d$ vanishes, since an integral coboundary in
the $T'$ argument has zero pairing. Integral Bloch--Kato
orthogonality for the de Rham representation $U$ identifies
the annihilator of all these tests with $H_{f,D}(T)$.
This proves $p^ef(y)\in H_{f,D}(T)$ on the dual side as well.
% END T84L3

% BEGIN T84L4
\paragraph{T84L4. The rank term and the full cut span.}
Let $A=H^1_f(K,T')$, the inverse image of the rational finite
condition in integral cohomology. Local finiteness makes $A$
finitely generated over $O$. Its rank is $\dim_L H^1_f(K,U')$:
a continuous cocycle on the compact group $G_K$ has bounded
image, so some nonzero scalar makes it integral. The same
denominator clearing applies to coboundaries. Over
$K=\mathbb Q_p$ the Bloch--Kato dimension formula gives
\[
 \dim_L H^1_f(K,U')
 =\dim_L (U')^{G_K}+\dim_L D_{\rm dR}(U')
                         -\dim_L {\rm Fil}^0D_{\rm dR}(U')\le k .
\]
Indeed $v\mapsto v\otimes1$ embeds the invariants in the
zero filtration, and $\dim_L D_{\rm dR}(U')=k$ by de Rhamness
and the degree-one hypothesis.

The torsion subgroup of $A$ needs at most $k$ generators too.
It lies in the kernel of $H^1(K,T')\to H^1(K,U')$, the image
of $(U'/T')^{G_K}$ by the connecting map. Take preimages of
finitely many generators of the torsion in this image. Their
span is a finite submodule of $(L/O)^k$, since their
denominators are jointly bounded. Such a submodule needs at
most $k$ generators by elementary divisors over a DVR. The
images of its generators generate the given torsion subgroup.
Consequently
\[
   \operatorname{length}_O(A/p^e A)\le
                2ke\operatorname{length}_O(O/pO).
\]
This also covers $e=0$ and $k=0$.

For the entire $O$-span $S$ of cut cocycles, Tate evaluation
now sends $S/(S\cap H_{f,D}(T))$ injectively into
$\operatorname{Hom}_O(A/p^e A,L/O)$. The zero kernel is precisely
Bloch--Kato orthogonality. The factorization through $A/p^e A$
holds for every generator by T84L3 and hence for the span.
The $O$-linear invariant map for $O(1)$ gives the scalar target
$L/O$ (equivalently choose a basis for the inverse different
in the scalar trace pairing). A finite $O$-module and its
$L/O$-dual have equal length, as one checks for $O/(\varpi^b)$
and then sums. The quotient on the left is
$(S+H_{f,D}(T))/H_{f,D}(T)$. This proves
$\mathrm{LocBoundP}(2ke)$ simultaneously for all test
lattices, without adding one cost per tuple or per point.
In particular both $\mathrm{StepLocDirect}$ and
$\mathrm{StepLocDual}$ of the type-prime cut follow from their
respective finite-cocycle lift and the assumed local
Fontaine, Bloch--Kato and Tate comparisons. When T84L1
provides that lift, $e=e_{\max}$. The production of the common
potential family and of the appropriate signed jump family,
and the quantitative hypothesis $\mathrm{admT}(C_T)$ used in
T84L1, are separate inputs.
% END T84L4

% BEGIN T84L5
\paragraph{T84L5. The two signs of the host jump condition.}
At $w\mid p$ put $c_0=(t-1)/2$, $H=s+S$, $a=\eta+t+1$
with $a<s$. The host test bounds $|U_w|\le n_V$ and contains
the jumps of every direct member at every embedding. Its
integer entries $b$ either lie between $0$ and $t-1$, satisfy
$|b-c_0|\le\eta+1$, or lie in the specified one-sided interval.
For the positive sign this last interval is
$c_0<b\le H$. All three alternatives lie in $[-a,H]$ since
$0\le c_0\le t-1$ and $a<s$. Thus $J=U_w\cup\{0\}$
suffices for every retained point.

For the negative sign the one-sided interval is
$2c_0-H\le b<c_0$. Reflect by $b\mapsto-b-1$. On this
interval the reflected value is at most $H-2c_0-1\le H$
and at least $-c_0-1\ge-a$. For $0\le b<t$ the reflected
value is in $[-t,-1]$. On the central interval it lies in
$[-c_0-\eta-2,\eta-c_0]$, again contained in $[-a,H]$
since $c_0+1\le t$. Take $J=\{-b-1:b\in U_w\}\cup\{0\}$.
Both chosen sets have size at most $n_V+1$, within the native
$N+1$ allowance. For an actual direct potential family each
point is de Rham by descent from $K'$. The dual matrix at $g$
is $\epsilon(g)R(g^{-1})^{\mathsf t}$ for $R=r_x\chi_1^{-1}$.
The dot-product pairing is perfect with values in $L(1)$;
the dual filtration reverses each labelled grading and $L(1)$
translates it by $-1$. Therefore a dual jump at $b'$ corresponds
exactly to a direct jump at $-b'-1$. This proves the two
$\mathrm{JumpFam}$ conclusions from the one host test.

The very same $K',F$ supply the dual potential family. Over $K'$
the Tate dual is crystalline. On crystalline periods the
cyclotomic line has basis $e_1\otimes t_{\rm cris}^{-1}$,
fixed under descent since the two cyclotomic factors cancel.
By the crystalline tensor and dual comparison, the pairing
on periods is perfect over $K'_0\otimes_{\mathbb Q_p}L$ and its
inertia action on that line is trivial (inertia fixes $K'_0$).
A component dual to the $\alpha$-space therefore transforms by
$\alpha^{-1}$. Inverting all characters preserves the
trivial character, their number, their restrictions over
$K'$, and their prime-to-$p$ orders; inversion commutes with
Frobenius conjugation. Thus $\Sigma^{-1}$ supplies
$\mathrm{PotFam}$ on the dual side without a new field or a
new size bound. Direct and dual matrix functions are continuous
by continuity of $r_x$, the characters, inversion in $G_K$ and
transpose. This supplies the sign-specific hypotheses of T84L1
from the direct potential family and the stated host test under
the native period comparisons.
% END T84L5

% BEGIN T84L6
\paragraph{T84L6. The rank block and each cyclotomic line at a type place.}
Let $\pi$ belong to the retained set, and use the given rational
semisimple decomposition of $r_\pi$ into the rank-$r$ block
$\rho_\pi$ and the characters $\chi_1,\ldots,\chi_t$. Twist
the same isomorphism by $\chi_1^{-1}$. Its inverse on the first
factor embeds $W_{\pi,O}$ rationally; on the $(c-1)$st
zero-based character coordinate it embeds
$\chi_c\chi_1^{-1}$ for every $1\le c\le t$. Each entry
of $\rho_\pi(g)$ is a fixed rational linear combination of the
entries of $r_\pi(g)$. Clearing its denominator gives an
integral continuous function; multiplication by a nonzero
element of a DVR is a homeomorphism onto its image in the
adic topology. The rank block therefore is continuous, as are
its twist and the character ratios. These embeddings give the
full $\mathrm{StepSub}$ tests for T84L2.

Set $q=c-1$. The original character formula on this range gives
$\chi_c\chi_1^{-1}=\epsilon^{-q}$, and the local coefficient
bridge identifies $\epsilon$ with the period cyclotomic
character. Under a coordinate basis the local lattice is thus
$O(-q)\subset L(-q)$. For the odd type prime $p$ and its
degree-one completion, if $q>0$ then the rational invariants
vanish (the cyclotomic image is infinite), and the de Rham
filtration has its unique step at $q>0$, so its degree-zero
quotient $D_{\rm dR}/{\rm Fil}^0$ is zero. The dimension formula
then gives $H^1_f(K,L(-q))=0$. Thus every finite rational
cocycle is a coboundary, and after reduction modulo $O(-q)$
the native finite cocycles are precisely the coboundaries.

For $q=0$ the same formula gives dimension one. The unramified
degree homomorphism is nonzero and crystalline: in the unramified
Witt periods solve $\varphi b-b=1$ by first solving
$z^p-z=1$ over the residue algebraic closure, then correcting
the error successively modulo higher powers of $p$ by the
same surjectivity of $z\mapsto z^p-z$. Inertia fixes this
period and $\varphi^m b-b=m$ on integral powers of arithmetic
Frobenius, hence on their closure by continuity. This spans
the rational finite condition. Reduction modulo $O$ therefore
vanishes on inertia. Conversely an unramified $L/O$-valued
homomorphism lifts on Frobenius to some $d\in L$, and the
continuous map $\widehat{\mathbb Z}\to\mathbb Z_p\to L$,
the last map being multiplication by $d$, gives the lift.
The native inertia identification gives the claimed unramified
test on the specified local bridge. This proves
$\mathrm{StepHfChar}$ and, with T84L1--T84L5, supplies the
rank-$r$ and all rank-one local estimates with exponents
$2r e_{\max}$ and $2 e_{\max}$ at every tested type place
once the direct potential family and $\mathrm{admT}(C_T)$
are given with the host, size and comparison hypotheses.
% END T84L6
